\section{Introduction}
\label{sec:intro}

Choosing a neural architecture means choosing how to spend a fixed budget of parameters and compute. The same
budget can buy depth (layers in sequence), breadth (parallel branches), or modularity
(sub-networks that each handle part of a problem). Modern systems increasingly mix all three of these types: multi-branch
blocks~\cite{szegedy2015going,xie2017resnext}, shared trunks with task heads~\cite{caruana1997multitask}, and
mixtures of experts~\cite{shazeer2017outrageously}. Theory says depth and width trade off in representational
power~\cite{eldan2016power,lu2017expressive}. People working on this, however, also care about a different concept:
\emph{how fast does the model answer?}

Latency is notoriously hard to predict from the numbers that architecture papers usually report. Parameter
count and FLOPs can disagree with the measured speed and can even reverse model
rankings~\cite{ma2018shufflenetv2,dehghani2022efficiency}. The gap is widest for small models, such as those in
embedded control, sensor fusion, or TinyML~\cite{banbury2021mlperftiny}, where fixed per-operator overhead,
rather than arithmetic, can dominate run time~\cite{sze2017efficient}. Benchmarks of such models are also easy to
get wrong. Framework choice, seed variance, and evaluation shortcuts can each produce differences larger
than the effect being studied~\cite{shi2016benchmarking,bouthillier2021accounting}.

This paper asks: \emph{at a fixed parameter budget, does the arrangement of parameters (sequential, parallel,
modular, or a hybrid) change accuracy and inference latency, and what predicts the latency?} We answer with a
controlled benchmark of six architectures at 385 parameters ($\pm$0.5\%) on three synthetic regression tasks of
increasing structural complexity. Synthetic tasks let us know the irreducible error exactly and let us design one task
whose compositional structure a modular model could benefit from.

\textbf{Contributions.}
\begin{itemize}
\item A \emph{parameter-matched benchmark} of three single-paradigm and three hybrid MLP architectures,
with five seeds, held-out evaluation, and a null control (two functionally identical models), released
with all code and raw results.
\item An \emph{audited pilot study} showing how a cross-system comparison produced a false
\pilotSpeedup$\times$ ``modular speedup'', and three low-cost practices that catch such errors.
\item Evidence that \emph{structure helps only when the task has structure}. All topologies tie on simple
targets, but a staged modular hybrid reduces error on a composed target by \chainHBGain\%, a gap that persists after
300 epochs of training.
\item A latency analysis showing that, for small models, \emph{single-sample latency is predicted by operator
count} ($R^2=\opsFitRsq$) rather than parameters or MACs, that \texttt{torch.compile} narrows but does not remove these differences,
and that the latency ranking \emph{depends on batch size}.
\end{itemize}

\Cref{sec:related} reviews related work, \cref{sec:method} describes the setup, \cref{sec:pilot} audits the pilot,
\cref{sec:results} presents results, and \cref{sec:discussion} discusses implications and limitations.
